\subsubsection{The Wave Equation}\label{subsec:wave-equation}

Our next elementary equation is the \emph{wave equation}. Here we seek a real-valued function $u$ depending on spatial variables $\bm{x}\in\mathbb{R}^n$ and a time variable $t\in\mathbb{R}$ satisfying
\begin{equation}\label{eq:wave-general}
u_{tt}=\Delta u.
\end{equation}
Once again the Laplacian acts only on the spatial variables. This equation describes many types of elastic and electromagnetic waves.

We once again describe some typical boundary conditions on the space-time cylinder $\Omega_t^+:=\{(\bm{x},t)\in\Omega\times(0,\infty)\}$, where $\Omega$ is a bounded domain in $\mathbb{R}^n$. Since the wave equation is second order in time one usually specifies two initial conditions
\begin{equation}\label{eq:wave-general-initial-position}
u(\bm{x},0)=f(\bm{x}),
\end{equation}

\begin{equation}\label{eq:wave-general-initial-velocity}
u_t(\bm{x},0)=g(\bm{x}).
\end{equation}
In problems in elasticity this amounts to specifying the position and velocity at time zero. Dirichlet or Neumann conditions are usually prescribed on various parts of the boundary. In elasticity applications these are usually interpreted as displacement and traction conditions, respectively.

\paragraph*{Solution of an initial/boundary-value problem by separation of variables}

The first initial/boundary-value problem we consider describes a string of unit length fixed at each end and given an initial position and velocity. The problem is described as follows. Let $D^+$ be the $(x,t)$ domain defined in the previous subsection. We seek a function $u:\overline{D^+}\to\mathbb{R}$ satisfying the one-dimensional (in space) wave equation
\begin{equation}\label{eq:wave-one-dimensional}
u_{tt}=u_{xx}
\end{equation}
for $(x,t)\in D^+$, the initial conditions
\begin{equation}\label{eq:wave-interval-initial-position}
u(x,0)=f(x),
\end{equation}

\begin{equation}\label{eq:wave-interval-initial-velocity}
u_t(x,0)=g(x)
\end{equation}
for $x\in(0,1)$, and the Dirichlet boundary conditions
\begin{equation}\label{eq:wave-left-boundary}
u(0,t)=0,
\end{equation}

\begin{equation}\label{eq:wave-right-boundary}
u(1,t)=0
\end{equation}
for $t>0$.

If we carry out the method of separation as before, we get the following family of solutions to both the wave equation \eqref{eq:wave-one-dimensional} and the boundary conditions.
\begin{equation}\label{eq:wave-separated-modes}
u_n(x,t)=(\alpha_n\cos n\pi t+\beta_n\sin n\pi t)\sin n\pi x.
\end{equation}


If our initial conditions have Fourier expansions of the form
\begin{equation}\label{eq:wave-position-fourier-series}
f(x)=\sum_{n=1}^{\infty}A_n\sin n\pi x,
\end{equation}

\begin{equation}\label{eq:wave-velocity-fourier-series}
g(x)=\sum_{n=1}^{\infty}B_n\sin n\pi x,
\end{equation}
then the formal series solution for the initial/boundary-value problem is
\begin{equation}\label{eq:wave-interval-series-solution}
u(x,t):=\sum_{n=1}^{\infty}\left(A_n\cos n\pi t+\dfrac{B_n}{n\pi}\sin n\pi t\right)\sin n\pi x.
\end{equation}


\paragraph*{D'Alembert's solution for the Cauchy problem}

In this section we consider the \emph{Cauchy problem} for the one-dimensional wave equation. Specifically, we wish to find a real-valued function $u$ that satisfies the wave equation \eqref{eq:wave-one-dimensional} in the half-plane $(x,t)\in(-\infty,\infty)\times(0,\infty)$ and the initial conditions
\begin{equation}\label{eq:wave-cauchy-initial-position}
u(x,0)=f(x),
\end{equation}

\begin{equation}\label{eq:wave-cauchy-initial-velocity}
u_t(x,0)=g(x)
\end{equation}
for $x\in(-\infty,\infty)$.

To derive a solution to this problem we first examine two special \emph{traveling wave} solutions of the wave equation. Suppose $F$ and $G$ are real-valued functions in $C^2(\mathbb{R})$. We observe that
\begin{equation}\label{eq:left-traveling-wave}
u_1(x,t):=F(x+t)
\end{equation}
and
\begin{equation}\label{eq:right-traveling-wave}
u_2(x,t):=G(x-t)
\end{equation}
each solve the wave equation. Note that $u_1$ is simply a translation of the function $F$ to the left with speed one, whereas $u_2$ is a translation of $G$ to the right. In fact, we can show that any solution of the wave equation has the form
\begin{equation}\label{eq:wave-traveling-sum}
u(x,t)=F(x+t)+G(x-t).
\end{equation}


To see this we simply make the change of variables
\begin{equation}\label{eq:wave-characteristic-xi}
\xi=x+t,
\end{equation}

\begin{equation}\label{eq:wave-characteristic-tau}
\tau=x-t,
\end{equation}
so that
\begin{equation}\label{eq:wave-transformed-function}
\bar u(\xi,\tau):=u\left(\dfrac{\xi+\tau}{2},\dfrac{\xi-\tau}{2}\right).
\end{equation}


Using the chain rule we see that if $u$ satisfies the wave equation then $\bar u$ satisfies
\begin{equation}\label{eq:wave-mixed-derivative-zero}
\bar u_{\xi\tau}=0.
\end{equation}
This implies
\begin{equation}\label{eq:wave-characteristic-split}
\bar u(\xi,\tau)=F(\xi)+G(\tau).
\end{equation}
Changing back to the independent variables $(x,t)$ gives us \eqref{eq:wave-traveling-sum}.

We now apply this general form for solutions to the Cauchy problem by plugging in the initial conditions \eqref{eq:wave-cauchy-initial-position} and \eqref{eq:wave-cauchy-initial-velocity} to get the following equations for the unknown functions $F$ and $G$:
\begin{equation}\label{eq:dalembert-initial-position-relation}
f(x)=F(x)+G(x),
\end{equation}

\begin{equation}\label{eq:dalembert-initial-velocity-relation}
g(x)=F'(x)-G'(x).
\end{equation}


These yield
\begin{equation}\label{eq:dalembert-f-derivative}
F'(x)=\dfrac12(f'(x)+g(x)),
\end{equation}

\begin{equation}\label{eq:dalembert-g-derivative}
G'(x)=\dfrac12(f'(x)-g(x)).
\end{equation}


Integrating these equations and using the result in \eqref{eq:wave-traveling-sum} gives us \emph{D'Alembert's solution of the Cauchy problem}
\begin{equation}\label{eq:dalembert-cauchy-solution}
u(x,t):=\dfrac12[f(x+t)+f(x-t)]+\dfrac12\int_{x-t}^{x+t}g(s)\,\mathrm{d}s.
\end{equation}


One of the most striking things about D'Alembert's solution (or more specifically, the form of the solution implied by \eqref{eq:wave-characteristic-split}) is that the formula for the solution makes perfectly good sense even when $f$ and $g$ are discontinuous. Such a "solution" would consist of a "jump" in $u$ moving to the left or right with unit speed. The existence of such solutions should not violate our intuition about the wave equation since physical wave-like phenomena that we would call discontinuous (such as breaking waves in the surf and shock waves from explosions) occur every day. But what about the mathematical nature of the solution? How can we say that a solution satisfies a differential equation at a point at which it is not differentiable? In later chapters we will examine this question more fully, and especially in the context of generalized wave equations we will get some fairly detailed answers.

\paragraph*{Energy conservation}

In this section we derive a result for solutions of the wave equation known as conservation of energy. We prove a version here that holds for the one-dimensional wave equation with fixed ends defined above and leave generalizations for later chapters.

\begin{Lemma}
\label{lem:wave-energy-conservation}
Let $u:\overline{D^+}\to\mathbb{R}$ be a $C^2$ solution of the wave equation \eqref{eq:wave-one-dimensional} satisfying the boundary conditions \eqref{eq:wave-left-boundary} and \eqref{eq:wave-right-boundary}. Then for any $t_1\geq t_0\geq0$, the solution $u$ satisfies
\begin{equation}\label{eq:wave-energy-conservation}
\int_0^1u_t^2(x,t_1)+u_x^2(x,t_1)\,\mathrm{d}x=\int_0^1u_t^2(x,t_0)+u_x^2(x,t_0)\,\mathrm{d}x.
\end{equation}

\end{Lemma}

\begin{Proof}
As we did in the proof of the energy inequality for the heat equation, we begin by deriving a differential identity. Let $u$ satisfy the wave equation. Then
\begin{equation}\label{eq:wave-equation-tested-by-velocity}
0=(u_{tt}-u_{xx})u_t
\end{equation}

\begin{equation}\label{eq:wave-local-energy-identity}
{}=\dfrac{\mathrm{d}}{\mathrm{d}t}(u_t^2+u_x^2)/2-\dfrac{\mathrm{d}}{\mathrm{d}x}(u_xu_t).
\end{equation}
We now use this in an integration over the rectangle $(x,t)\in[0,1]\times[t_0,t_1]$, in which we change the order of integration at will, and we obtain the following:
\[
\begin{aligned}
&\int_0^1u_t^2(x,t_1)+u_x^2(x,t_1)\,\mathrm{d}x
-\int_0^1u_t^2(x,t_0)+u_x^2(x,t_0)\,\mathrm{d}x\\
&\quad=\int_0^1\int_{t_0}^{t_1}\dfrac{\mathrm{d}}{\mathrm{d}t}(u_t^2+u_x^2)\,\mathrm{d}t\,\mathrm{d}x\\
&\quad=\int_{t_0}^{t_1}\int_0^12\dfrac{\mathrm{d}}{\mathrm{d}x}(u_xu_t)\,\mathrm{d}x\,\mathrm{d}t\\
&\quad=\int_{t_0}^{t_1}2u_x(1,t)u_t(1,t)\,\mathrm{d}t-\int_{t_0}^{t_1}2u_x(0,t)u_t(0,t)\,\mathrm{d}t.
\end{aligned}
\]
However, the boundary conditions \eqref{eq:wave-left-boundary} and \eqref{eq:wave-right-boundary} imply
\begin{equation}\label{eq:wave-boundary-time-derivatives}
u_t(0,t)=u_t(1,t)=0,
\end{equation}
so this gives us \eqref{eq:wave-energy-conservation}.
\end{Proof}

Note that the quantity we call the energy for solutions of the wave equation and the quantity we call the energy for solutions of the heat equation seem very different mathematically. However, the mathematical techniques that we use to study the quantities (multiplication of the differential equation by the solution or its derivative and (essentially) integrating by parts in order to obtain an estimate) are common to both. This technique of obtaining estimates on solutions of PDEs is extremely useful and is generalized in later chapters.

\paragraph*{Problems}

\begin{enumerate}[label=\thesection.\arabic*.,ref=\thesection.\arabic*,start=25,series=book-problems]

\item\label{ex:wave-two-boundary-problems} Solve the one-dimensional wave equation via separation of variables for the following boundary conditions:\begin{enumerate}[label=(\alph*)]
\item $\begin{aligned}u(x,0)&=x^2-x,\\u_t(x,0)&=0,\\u_x(0,t)&=0,\\u_x(1,t)&=0.\end{aligned}$
\item $\begin{aligned}u(x,0)&=0,\\u_t(x,0)&=\sin\pi x,\\u(0,t)&=0,\\u_x(1,t)&=0.\end{aligned}$
\end{enumerate}

\item\label{ex:wave-cauchy-wellposedness} Give a specific definition of well-posedness (in particular, make precise in what sense the problem is continuous with respect to the data) for the Cauchy problem \eqref{eq:wave-one-dimensional}, \eqref{eq:wave-cauchy-initial-position}, \eqref{eq:wave-cauchy-initial-velocity} on the domain $(x,t)\in(-\infty,\infty)\times(0,\infty)$. Derive conditions on the initial data under which the problem is well-posed. How do your results differ if the domain under consideration is $(x,t)\in(-\infty,\infty)\times(0,T)$ for some $0<T<\infty$. Hint: If $u(x,0)=0$ and $u_t(x,0)=\epsilon>0$ for $x\in(-\infty,\infty)$, then $u$ grows arbitrarily large with time. Figure out conditions on the initial data that assure that $u$ stays bounded.

\item\label{ex:wave-finite-propagation-region} Suppose $f$ and $g$ are identically zero outside the interval $[-1,1]$. In what region in $(-\infty,\infty)\times[0,\infty)$ can you ensure that the solution $u$ of the Cauchy problem is identically zero.

\item\label{ex:heat-propagation-comparison} Is there a similar result to the previous problem for the heat equation? Hint: Use
\[
\chi_{[-1,1]}(x):=\begin{cases}1,&x\in[-1,1]\\0,&x\notin[-1,1]\end{cases}
\]
as initial datum. Use Problem~\ref{ex:heat-step-cauchy-problem} to obtain a solution.

\item\label{ex:discontinuous-weak-wave-solutions} We define a weak solution of the one-dimensional wave equation to be a function $u(x,t)$ such that
\begin{equation}\label{eq:wave-distributional-form}
\int_{-\infty}^{\infty}\int_{-\infty}^{\infty}u(x,t)(\phi_{tt}(x,t)-\phi_{xx}(x,t))\,\mathrm{d}x\,\mathrm{d}t=0
\end{equation}
for every $\phi\in C_0^2(\mathbb{R}^2)$. Here $C_0^2(\mathbb{R}^2)$ is the set of functions in $C^2(\mathbb{R}^2)$ that have compact support; i.e., that are identically zero outside of some bounded set.
\begin{enumerate}[label=(\alph*)]
\item Show that any strong (classical $C^2$) solution of the wave equation is also a weak solution.
\item Show that discontinuous functions of the form
\begin{equation}\label{eq:heaviside-right-moving-solution}
u(x,t):=H(x-t)
\end{equation}
and
\begin{equation}\label{eq:heaviside-left-moving-solution}
u(x,t):=H(x+t)
\end{equation}
are weak solutions of the wave equation. Here $H$ is the Heaviside function:
\begin{equation}\label{eq:heaviside-function}
H(x):=\begin{cases}0,&x<0\\1,&x\geq0.\end{cases}
\end{equation}
\end{enumerate}

\end{enumerate}
